\begin{explanationbox}
\textbf{\textcolor{CExplanation}{同一个椭圆：极线存在，不等于有切点弦}}

\begin{center}
\begin{tikzpicture}[x=.85cm,y=.85cm,>=Stealth,font=\small]
 \path[use as bounding box] (-2.6,-2) rectangle (12.3,2.15);
 % Both panels use exactly the same coordinate scale and ellipse.
 \begin{scope}
  \node at (1,1.8) {外点：两个实切点};
  \draw[->,gray] (-2.3,0)--(4.5,0) node[right] {$x$};
  \draw[CStatement,thick] (0,0) ellipse (2 and 1);
  \draw[orange!80!black,thick] (1,.8660254)--(4,0)--(1,-.8660254);
  \draw[teal!80!black,very thick] (1,-1.25)--(1,1.25);
  \fill (4,0) circle(2pt) node[below] {$P(4,0)$};
  \fill (1,.8660254) circle(2pt); \fill (1,-.8660254) circle(2pt);
  \node[teal!80!black] at (.7,-1.6) {极线 $x=1$};
 \end{scope}
 \begin{scope}[shift={(7.5,0)}]
  \node at (1,1.8) {内点：没有实切点};
  \draw[->,gray] (-2.3,0)--(4.5,0) node[right] {$x$};
  \draw[CStatement,thick] (0,0) ellipse (2 and 1);
  \draw[teal!80!black,very thick] (4,-1.25)--(4,1.25);
  \fill (1,0) circle(2pt) node[below] {$P(1,0)$};
  \node[teal!80!black] at (3,-1.6) {极线 $x=4$};
 \end{scope}
\end{tikzpicture}

\end{center}

两图均为 $x^2/4+y^2=1$，且坐标尺度相同。
左图 $P(4,0)$ 的极线 $x=1$ 经过两个切点；
右图 $P(1,0)$ 的极线 $x=4$ 却在椭圆之外，因为椭圆上 $|x|\le2$。
\textbf{先检查是否成直线，再检查是否有两个不同实切点。}

\medskip
\textbf{\textcolor{CExplanation}{一句话理解}}
\textbf{先从切点写切线，再把“切点”和“给定点”的坐标交换位置。}
能交换，是因为公式中的 $x_0x$、$y_0y$、$x_0y+y_0x$ 本来就是左右交换后不变的组合。

\medskip
\begin{minipage}{\linewidth}
\textbf{\textcolor{CExplanation}{一、圆为什么能产生这种直线？}}
圆 $x^2+y^2=R^2$ 上一点 $T(u,v)$，半径 $OT$ 与过 $T$ 的切线垂直。
设切线上任一点为 $Q(x,y)$，则
\[
 \overrightarrow{OT}=(u,v),\qquad
 \overrightarrow{TQ}=(x-u,y-v),
\]
两向量垂直，故
\[
 u(x-u)+v(y-v)=0
 \ \Longrightarrow\ ux+vy=u^2+v^2=R^2.
\]
若从 $P(x_0,y_0)$ 引切线，切点 $T$ 就必须满足 $ux_0+vy_0=R^2$。
把 $T(u,v)$ 看作待求点 $(x,y)$，便得到 $x_0x+y_0y=R^2$：\textbf{它是所有切点的共同直线}。
椭圆可以将单位圆的横坐标乘 $a$、纵坐标乘 $b$ 得到，故圆的这个规律也变成椭圆公式。

\end{minipage}
\par
\medskip
\begin{minipage}{\linewidth}
\textbf{\textcolor{CExplanation}{二、为什么 $xy$、$x$、$y$ 项要除以 $2$？}}
设切点坐标为 $(u,v)$，一般二次曲线的切线中有下列对称组合：
\[
 x^2\longrightarrow ux,\quad
 y^2\longrightarrow vy,\quad
 xy\longrightarrow\frac{uy+vx}{2},\quad
 x\longrightarrow\frac{x+u}{2}.
\]
例如把 $Q$ 取成切点 $T$，就有
\[
 \left.\frac{uy+vx}{2}\right|_{(x,y)=(u,v)}=uv,
 \qquad
 \left.\frac{x+u}{2}\right|_{x=u}=u.
\]
因此，\textbf{在两个点重合时，每一项都还原成原曲线方程中的对应项}。
这既能检验公式，也能防止交叉项漏写一半；完整理由见第四种证明。

\end{minipage}
\par
\medskip
\begin{minipage}{\linewidth}
\textbf{\textcolor{CExplanation}{三、别把“有极线”误解成“有切点弦”}}
\begin{itemize}[leftmargin=*,itemsep=0.35em]
 \item \textbf{椭圆内点：}如本页右图，极线在椭圆之外，没有真实切点。
 \item \textbf{双曲线上的特殊外点：}$x^2-y^2=1$、$P(1,1)$ 只有一个实切点，见例 4 的构造图；不能叫两切点连线。
 \item \textbf{中心例外：}例如椭圆或双曲线的中心 $O(0,0)$，替换后得到 $0=1$，不是直线。
\end{itemize}

\end{minipage}
\par
\medskip
\begin{minipage}{\linewidth}
\textbf{\textcolor{CExplanation}{四、解题时怎么识别？}}
\begin{description}[style=nextline,leftmargin=2.8em,labelsep=0.5em,itemsep=0.35em]
 \item[\textbf{信号一：已知曲线上一点，求切线。}] 直接把该点的坐标写入极线公式；它就是切线。
 \item[\textbf{信号二：已知曲线外一点，求两个切点连线。}] 先写极线，再与原曲线联立，确认有两个不同的实切点。
 \item[\textbf{信号三：给出一般二次方程。}] 注意 $xy$、$x$、$y$ 三类项的“平均写法”，最后检验所得等式确为直线。
 \item[\textbf{信号四：求定点或轨迹。}] 把极线整理成 $Lx+Ny+K=0$，通过对比系数，寻找定点坐标之间的关系。
\end{description}
\end{minipage}
\par
\end{explanationbox}
